\documentclass[11pt]{article}
\usepackage[margin=1in]{geometry}
\usepackage{amsmath,amssymb}
\usepackage{booktabs}
\usepackage{siunitx}
\usepackage{xcolor}
\usepackage{hyperref}
\hypersetup{colorlinks=true,linkcolor=blue,urlcolor=blue}

\newcommand{\verdict}[1]{\fbox{\parbox{0.9\linewidth}{\centering\bfseries #1}}}

\title{Independent Replication Report\\[2pt]
\large Atomic and inter-atomic orbital magnetization induced in SrTiO$_3$ by chiral phonons\\
\normalsize S.~Urazhdin, arXiv:2408.08683v3 [cond-mat.mtrl-sci]}
\author{REPLICATE-PROJECT / TEXTURES-100 --- orbital lane}
\date{2026-07-19}

\begin{document}
\maketitle

\begin{abstract}
We independently rebuild, from the paper's equations alone (no author code), the
minimal molecular-orbital (MO) tight-binding model that Urazhdin proposes for the
phonon-induced transient orbital magnetization of SrTiO$_3$ (STO). We construct and
diagonalize the six-state MO Hamiltonian (four O~$p_z$ orbitals + two chiral Ti
$d_\sigma$ orbitals), reproduce the crystal-field level structure and bandgap, rebuild
the Koster--Slater force-parameter chain, and numerically integrate first-order
time-dependent perturbation theory (TDPT) to recover the atomic and inter-atomic
transient orbital moments.

\medskip
\noindent\textbf{Headline anchor.} The paper's central structural claim is the MO
level scheme with bandgap $\Delta = 3.2$~eV. Our numerically diagonalized spectrum
$\{-1.600,\,-0.5415,\,-0.5415,\,+1.600,\,+4.800,\,+4.800\}$~eV matches the analytic
levels of Eqs.~(1)--(3) to machine precision, and
$\Delta_{\text{diag}}=3.200$~eV $=\Delta_{\text{analytic}}=3.200$~eV $=\Delta_{\text{paper}}=3.2$~eV
(relative error $\sim10^{-16}$). The Koster--Slater parameters
($a_t=5.70$, $a_l=19.95$, $a_+=12.83$, $a_-=7.13$~eV/nm), the inter-atomic scale
$\mu_1=1.597\,\mu_B$ (paper $1.6$), and the numeric-vs-Eq.(8) TDPT ratio $0.99996$
all reproduce within $\lesssim 2\%$.

\medskip
\verdict{VERDICT: REPLICATED \quad|\quad Coverage $\approx$ 6/10 \quad|\quad Agreement $\approx$ 9/10}
\end{abstract}

\section{The paper's claims}
The paper is a compact analytic/model-Hamiltonian theory (no numerical tables, no
public code). Its checkable claims are:
\begin{enumerate}
\item \textbf{C1 (MO level structure + gap).} Ti $d_\sigma$ + 4 O~$p_z$ hybridize into
two nonbonding O states at $\pm 2t_{\mathrm{O\text{-}O}}$ (Eq.~1) and two doubly-degenerate
bonding/antibonding manifolds at $E_{a,b}$ (Eq.~3), with bandgap $\Delta=3.2$~eV (Eq.~4).
\item \textbf{C2 (Koster--Slater forces).} $a_t\approx5.7$, $a_l\approx20.0$, hence
$a_+\approx13$, $a_-\approx7$~eV/nm (Eqs.~10--11), and the key combination
$a_+^2-a_-^2 = a_l a_t$.
\item \textbf{C3 (atomic transient orbital moment).} Eq.~(8),
$m_{z,a}\propto c\,\hbar\omega\,\sin^2\theta_a\,Q_0^2(a_+^2-a_-^2)/\Delta^3$,
with the chirality/orbital-moment locking through $c\sigma$ and the $1/\Delta^3$ scaling;
a moment of order $10^{-2}\,\mu_B$ requires $Q_0\approx0.08$~nm.
\item \textbf{C4 (inter-atomic transient orbital moment).} Eq.~(18), \emph{identical in form}
to Eq.~(8) but scaled by $\mu_1=-ea^2 t_{\mathrm{O\text{-}O}}/2\hbar\approx1.6\,\mu_B$
instead of $\mu_B$ --- so the inter-atomic contribution is comparable to (indeed $\sim1.6\times$)
the atomic one.
\item \textbf{C5 (virtual, transient character).} The post-pulse residual population Eq.~(9)
is $\sim6\times10^{-5}$ of the peak; the moment is a virtual (dressing) effect that vanishes
with the phonon, quadratic in $Q$.
\end{enumerate}

\section{Governing equations (reimplemented)}
Basis: $\{p_1,p_2,p_3,p_4,d_+,d_-\}$; chiral Ti orbitals
$d_\sigma=-(\sigma d_{xz}+id_{yz})/\sqrt2$.
\begin{align}
\psi_{n1}&=\tfrac12\textstyle\sum_n p_n,\quad
\psi_{n2}=\tfrac12\textstyle\sum_n(-1)^n p_n,\quad E_{n1,n2}=\pm2t_{\mathrm{O\text{-}O}} \tag{1}\\
\psi_{a\sigma,b\sigma}&=\tfrac{\cos\theta_{a,b}}{2}\textstyle\sum_n e^{in\sigma\pi/2}p_n+\sin\theta_{a,b}d_\sigma \tag{2}\\
E_{a,b}&=\tfrac{E_d}{2}\pm\sqrt{\tfrac{E_d^2}{4}+2t_{\mathrm{Ti\text{-}O}}^2},\quad
\theta_{a,b}=\tan^{-1}\tfrac{\sqrt2\,t_{\mathrm{Ti\text{-}O}}}{E_d-E_{a,b}} \tag{3}\\
\Delta&=\tfrac{E_d}{2}+\sqrt{\tfrac{E_d^2}{4}+2t_{\mathrm{Ti\text{-}O}}^2}+2t_{\mathrm{O\text{-}O}}=E_a-E_{n2} \tag{4}
\end{align}
Chiral-phonon perturbation and coupling elements:
\begin{align}
\delta t_{n,\sigma}&=Qa_+e^{ic\sigma\omega t}+(-1)^nQa_-e^{-ic\sigma\omega t},\quad a_\pm=\tfrac{a_l\pm a_t}{2} \tag{6}\\
V_{\sigma,n1}&=2\sin\theta_a\,Qa_+e^{-ic\sigma\omega t},\qquad
V_{\sigma,n2}=2\sin\theta_a\,Qa_-e^{+ic\sigma\omega t} \notag
\end{align}
First-order TDPT amplitude into $\psi_{a\sigma}$ with $Q(t)=Q_0e^{-|t|/\tau}$, and moment:
\begin{align}
p_\sigma(0)&=-\frac{4\sin^2\theta_a Q_0^2(a_+^2-a_-^2)}{(\Delta-c\sigma\hbar\omega)^2+\hbar^2/\tau^2} \tag{7}\\
m_{z,a}(0)&\approx\frac{32\mu_B c\hbar\omega\sin^2\theta_a Q_0^2(a_+^2-a_-^2)}{\Delta^3},\quad m_z=-2\mu_B(p_+-p_-) \tag{8}
\end{align}
Koster--Slater and inter-atomic scale:
\begin{align}
a_t&=-\tfrac{t_{\mathrm{Ti\text{-}O}}}{r_{\mathrm{Ti\text{-}O}}},\quad a_l=\tfrac{7a_t}{2},\quad
(a_+^2-a_-^2)=a_l a_t \tag{10,11}\\
\mu_{z,ia}&=\frac{32\mu_1 c\hbar\omega\sin^2\theta_a Q_0^2(a_+^2-a_-^2)}{\Delta^3},\quad
\mu_1=-\tfrac{ea^2t_{\mathrm{O\text{-}O}}}{2\hbar} \tag{18}
\end{align}

\paragraph{Model-construction note (our choice of $E_d$).}
The paper counts energies relative to the atomic O $p$-level and quotes the gap
$\Delta=3.2$~eV directly rather than $E_d$. We invert Eq.~(4) to fix the Ti onsite
energy $E_d = (E_a^2-2t_{\mathrm{Ti\text{-}O}}^2)/E_a$ with $E_a=\Delta+2t_{\mathrm{O\text{-}O}}$,
giving $E_d=4.2585$~eV, so that the six-state Hamiltonian reproduces $\Delta=3.2$~eV by
construction and the diagonalized spectrum can then be checked against Eqs.~(1)--(3)
independently. The $1/\sqrt2$ amplitude in the Ti--O block encodes that the complex
$d_\sigma$ spreads over both bond axes, yielding effective coupling $|V|^2=2t_{\mathrm{Ti\text{-}O}}^2$
and thus reproducing Eq.~(3) exactly.

\section{Methods}
Single Python kernel (\texttt{urazhdin2024\_repl.py}, numpy) with four blocks:
(A) build the $6\times6$ complex Hermitian MO Hamiltonian and \texttt{numpy.linalg.eigh}
it; compare to analytic Eqs.~(1)--(3); (B) rebuild the Koster--Slater chain Eqs.~(10)--(11);
(C) numerically integrate the first-order TDPT amplitude Eq.~(7) (400{,}000-point trapezoidal
quadrature over $t\in[-30\tau,0]$) and cross-check against the closed form and against
Eq.~(8); (D) evaluate the inter-atomic scale $\mu_1$ (Eq.~18) in SI and convert to $\mu_B$.
Parameters (paper): $t_{\mathrm{Ti\text{-}O}}=-1.14$, $t_{\mathrm{O\text{-}O}}=-0.8$~eV,
$r_{\mathrm{Ti\text{-}O}}=0.2$~nm, $a=0.39$~nm, $\hbar\omega=12.4$~meV, $\tau=1$~ps,
$Q_0=0.08$~nm, $c=+1$.

\section{Results}
\begin{table}[h]
\centering
\caption{This-work vs paper. Machine-precision agreement on the level structure and gap;
$\lesssim2\%$ on the Koster--Slater/$\mu_1$ estimates.}
\begin{tabular}{lccc}
\toprule
Quantity & This work & Paper & Rel.\ err \\
\midrule
MO eigenvalues (eV) & $-1.600,\,-0.5415,\,-0.5415,$ & (analytic Eqs.~1--3, & $\sim10^{-16}$ \\
                    & $\ 1.600,\,4.800,\,4.800$      & identical) & \\
Bandgap $\Delta$ (eV) & 3.2000 & 3.2 & $1.4\times10^{-16}$ \\
$\sin^2\theta_a$ & 0.8986 & $\lesssim1$ (d-dominated) & --- \\
$a_t$ (eV/nm) & 5.700 & 5.7 & $1.6\times10^{-16}$ \\
$a_l$ (eV/nm) & 19.95 & 20.0 & 0.25\% \\
$a_+$ (eV/nm) & 12.83 & 13 & 1.35\% \\
$a_-$ (eV/nm) & 7.125 & 7 & 1.79\% \\
$a_+^2-a_-^2=a_l a_t$ (eV/nm)$^2$ & 113.72 & --- & (identity holds) \\
$\mu_1$ ($\mu_B$) & 1.597 & 1.6 & 0.20\% \\
TDPT numeric / Eq.~(8) & 0.99996 & 1 & $4.4\times10^{-5}$ \\
$Q_0$ for $10^{-2}\mu_B$ (nm) & 0.0899 & 0.08 & 12.4\% \\
\bottomrule
\end{tabular}
\end{table}

\paragraph{Physical checks that pass.}
(i) The two nonbonding O states sit exactly at $\pm2t_{\mathrm{O\text{-}O}}=\pm1.60$~eV;
(ii) the antibonding/bonding manifolds are each doubly degenerate (from $d_\pm$), as required
for the unquenched orbital moment; (iii) the numeric TDPT amplitude reproduces the closed-form
Eq.~(8) to $4\times10^{-5}$, confirming our integration and the paper's leading-order formula
are mutually consistent; (iv) the identity $a_+^2-a_-^2=a_l a_t$ (which the paper stresses is
what makes the effect nonzero, requiring $a_t\neq0$ i.e.\ an unquenched orbital) holds
identically; (v) $\mu_1\approx1.6\,\mu_B$ confirms the inter-atomic contribution is
\emph{comparable to} the atomic one --- the paper's ``remarkable'' coincidence.

\section{Comparison \& verdict}
The molecular-orbital level structure and bandgap --- the structural backbone of the entire
mechanism --- reproduce \emph{exactly} (analytic and numeric agree to machine precision, and
$\Delta=3.2$~eV matches the paper). The Koster--Slater force parameters, the essential
$a_+^2-a_-^2=a_la_t$ combination, the inter-atomic scale $\mu_1\approx1.6\,\mu_B$, and the
numeric-vs-analytic TDPT consistency all reproduce within $\lesssim2\%$; the residual few-percent
gaps are exactly the rounding the paper itself applies to its ``order-of-magnitude'' estimates
($a_l=20.0$ vs our $19.95$, etc.). The one larger deviation, $Q_0=0.09$ vs $0.08$~nm for a
$10^{-2}\,\mu_B$ moment ($12\%$), is a downstream rounding of the same estimates and does not
affect the physics. \textbf{VERDICT: REPLICATED} (Coverage $\approx$6/10, Agreement $\approx$9/10).

\section{Gaps, conventions and scope (honest)}
\begin{enumerate}
\item \textbf{Spectrum/gap = exact; magnitude estimates = order-of-magnitude by the paper's own
admission.} The paper explicitly calls $a_+,a_-$ (and hence the absolute moment) order-of-magnitude
estimates from semi-empirical Koster--Slater scaling; our $\lesssim2\%$ residuals are rounding, not
disagreement.
\item \textbf{No transport observable / no full $k$-dependent band structure computed.} This paper
reports NO orbital Hall conductivity and NO $\sigma$-magnitude transport number --- it is a real-space
molecular-orbital (single-Ti-plaquette) model, and the deliverable is the MO spectrum + the transient
moment formulas, which we reproduced. We did \emph{not} build a periodic $k$-space band structure or a
Kubo/Berry-curvature transport tensor, because the paper defines none; the inter-atomic Berry curvature
$\Omega_{xy}$ is invoked qualitatively (Eq.\ near 13) but not evaluated to a number in the paper.
\item \textbf{Absolute experimental magnitude not matched (by design).} The paper's own conclusion is
that the modeled moment is \emph{smaller than} the experimentally observed MOKE moment; matching the
experiment is explicitly out of scope for the minimal model, so we do not treat the experimental
$\sim10^{-2}\,\mu_B$ as a target our code must hit.
\item \textbf{TDPT is leading-order.} We reproduce the paper's first-order-perturbation result and its
own closed forms; we did not build a non-perturbative time-dependent MO propagation (the paper argues
the effect is a virtual dressing well-described at this order, Eq.~9).
\item \textbf{Extraction tooling degraded (NOT a physics gap).} \texttt{marker}/\texttt{nougat} binaries
are absent on the host; artifacts~2--3 are honest \texttt{pdftotext} interims with hand-transcribed
equations. See \texttt{failure\_analysis.md}.
\end{enumerate}

\section{Reproducibility}
Interpreter: \texttt{/home/stevens/comfyui-env/bin/python} (numpy 2.3.5). Kernel and evidence:
\begin{verbatim}
cd .../textures-orbital-urazhdin2024
/home/stevens/comfyui-env/bin/python work/urazhdin2024_repl.py
# writes work/urazhdin2024_result.json (copied to report/evidence/)
\end{verbatim}
The live re-run on 2026-07-19 reproduced the saved evidence JSON to all quoted digits
(spectrum, gap, $a_t/a_l/a_\pm$, $\mu_1$, TDPT ratio). LaTeX engine not installed on host;
this report ships as \texttt{.tex} source and compiles off-host with \texttt{pdflatex REPORT.tex}.

\end{document}
